% ============================================================
%  Genomic Syntactic Analysis
%  A Formal Grammar Approach to Genome Structure
% ============================================================
\documentclass[12pt, a4paper, oneside]{book}

% ── Encoding & Fonts ────────────────────────────────────────
\usepackage[T1]{fontenc}
\usepackage[utf8]{inputenc}
% \usepackage{microtype}  % disabled: requires scalable fonts
\emergencystretch=3em   % allow extra stretch to eliminate overfull hboxes

% ── Mathematics ─────────────────────────────────────────────
\usepackage{amsmath, amssymb, amsthm}
\usepackage{mathtools}

% ── Layout ──────────────────────────────────────────────────
\usepackage[
  a4paper,
  top=30mm, bottom=30mm,
  left=35mm, right=25mm,
  headheight=15pt
]{geometry}
\usepackage{setspace}
\onehalfspacing

% ── Headers & Footers ───────────────────────────────────────
\usepackage{fancyhdr}
\pagestyle{fancy}
\fancyhf{}
\fancyhead[L]{\small\slshape\leftmark}
\fancyhead[R]{\small\thepage}
\renewcommand{\headrulewidth}{0.4pt}

% ── Colors ──────────────────────────────────────────────────
\usepackage[dvipsnames, table]{xcolor}
\definecolor{c1}{HTML}{0ea5e9}   % sky blue  – level accent
\definecolor{c2}{HTML}{7c3aed}   % violet
\definecolor{c3}{HTML}{10b981}   % emerald
\definecolor{c4}{HTML}{f59e0b}   % amber
\definecolor{c5}{HTML}{ef4444}   % red
\definecolor{c6}{HTML}{ec4899}   % pink
\definecolor{bgdark}{HTML}{0f172a}
\definecolor{bggray}{HTML}{f8fafc}
\definecolor{rulegray}{HTML}{e2e8f0}
\definecolor{ntcolor}{HTML}{10b981}
\definecolor{tcolor}{HTML}{f59e0b}
\definecolor{commentcolor}{HTML}{94a3b8}

% ── Graphics & TikZ ─────────────────────────────────────────
\usepackage{tikz}
\usetikzlibrary{
  positioning, arrows.meta, shapes.geometric,
  backgrounds, fit, decorations.pathreplacing,
  calc, shadows.blur
}
\usepackage{pgfplots}
\pgfplotsset{compat=1.18}

% ── Tables ──────────────────────────────────────────────────
\usepackage{booktabs}
\usepackage{tabularx}
\usepackage{multirow}
\usepackage{array}
\usepackage{colortbl}

% ── Code / Verbatim ─────────────────────────────────────────
\usepackage{listings}
\lstset{
  basicstyle=\ttfamily\small,
  backgroundcolor=\color{bggray},
  frame=single,
  framerule=0pt,
  rulecolor=\color{rulegray},
  breaklines=true,
  showstringspaces=false,
  commentstyle=\color{commentcolor}\itshape,
  keywordstyle=\color{c2}\bfseries,
  stringstyle=\color{c3},
  numbers=left,
  numberstyle=\tiny\color{commentcolor},
  xleftmargin=18pt,
}

% ── Boxes & Environments ────────────────────────────────────
\usepackage[most]{tcolorbox}
\tcbuselibrary{skins, breakable, listings}

\newtcolorbox{grammarbox}[1][]{
  enhanced,
  colback=bggray,
  colframe=c2,
  arc=4pt,
  boxrule=0.6pt,
  left=10pt, right=10pt, top=8pt, bottom=8pt,
  fonttitle=\bfseries\small,
  title={Generative Grammar Rules},
  attach boxed title to top left={yshift=-2mm, xshift=6mm},
  boxed title style={colback=c2, arc=3pt, boxrule=0pt},
  #1
}

\newtcolorbox{notebox}[1][Note]{
  enhanced,
  colback=c1!6!white,
  colframe=c1,
  arc=4pt,
  boxrule=0.5pt,
  left=8pt, right=8pt, top=6pt, bottom=6pt,
  fonttitle=\bfseries\small,
  title={#1},
  attach boxed title to top left={yshift=-2mm, xshift=6mm},
  boxed title style={colback=c1, arc=3pt, boxrule=0pt},
}

\newtcolorbox{keybox}[1][Key Concept]{
  enhanced,
  colback=c3!6!white,
  colframe=c3,
  arc=4pt,
  boxrule=0.5pt,
  left=8pt, right=8pt, top=6pt, bottom=6pt,
  fonttitle=\bfseries\small,
  title={#1},
  attach boxed title to top left={yshift=-2mm, xshift=6mm},
  boxed title style={colback=c3, arc=3pt, boxrule=0pt},
}

% ── Hyperlinks (load last) ───────────────────────────────────
\usepackage[
  colorlinks=true,
  linkcolor=c2,
  urlcolor=c1,
  citecolor=c3,
  bookmarksnumbered=true,
  pdfusetitle
]{hyperref}
\usepackage{cleveref}

% ── Theorem Environments ────────────────────────────────────
\theoremstyle{definition}
\newtheorem{definition}{Definition}[chapter]
\newtheorem{example}{Example}[chapter]
\theoremstyle{remark}
\newtheorem{remark}{Remark}[chapter]

% ── Custom Commands ─────────────────────────────────────────
\newcommand{\NT}[1]{\textcolor{ntcolor}{\texttt{#1}}}   % Non-terminal
\newcommand{\T}[1]{\textcolor{tcolor}{\texttt{#1}}}     % Terminal
\newcommand{\derives}{\;\longrightarrow\;}
\newcommand{\gor}{\;\mid\;}                              % grammar OR
\newcommand{\bp}{\,\text{bp}}
\newcommand{\kbp}{\,\text{kbp}}
\newcommand{\Mbp}{\,\text{Mbp}}
\newcommand{\Gbp}{\,\text{Gbp}}

% ── Document Metadata ───────────────────────────────────────
\title{Genomic Syntactic Analysis:\\
       A Formal Grammar Approach to\\
       Genome Structure and Variation}
\author{Wanhong\\[4pt]
        \small Serendipitech Inc.}
\date{\today}

% ============================================================
\begin{document}
% ============================================================

% ── TITLE PAGE ──────────────────────────────────────────────
\begin{titlepage}
\pagecolor{bgdark}
\color{white}

\begin{tikzpicture}[remember picture, overlay]
  % Background decorative arcs
  \draw[c2!40, line width=0.4pt, opacity=0.5]
    (current page.center) ++(-6,4) arc(120:420:8cm);
  \draw[c1!40, line width=0.3pt, opacity=0.4]
    (current page.center) ++(5,-3) arc(0:270:11cm);
  % DNA helix hint (stylized dots)
  \foreach \y in {0,0.5,...,10}{
    \pgfmathsetmacro{\xA}{0.5*cos(\y*60)}
    \pgfmathsetmacro{\xB}{-0.5*cos(\y*60)}
    \fill[c3!60, opacity=0.35]
      ([xshift=\xA cm, yshift=\y cm-1cm] current page.south west)
      ++( 2.2,0) circle (1.5pt);
    \fill[c4!60, opacity=0.35]
      ([xshift=\xB cm, yshift=\y cm-1cm] current page.south west)
      ++( 2.2,0) circle (1.5pt);
  }
\end{tikzpicture}

\vspace*{3cm}

\begin{center}

{\fontsize{9}{11}\selectfont\ttfamily\color{c1}\MakeUppercase{%
  Research Monograph \quad|\quad Computational Genomics Series}}

\vspace{0.8cm}
\rule{\textwidth}{0.5pt}
\vspace{0.6cm}

{\fontsize{32}{38}\selectfont\bfseries
Genomic Syntactic Analysis}

\vspace{0.4cm}

{\fontsize{16}{20}\selectfont\color{c1}
A Formal Grammar Approach to\\[4pt]
Genome Structure and Variation}

\vspace{0.6cm}
\rule{\textwidth}{0.5pt}
\vspace{1.5cm}

{\large Wanhong}\\[6pt]
{\normalsize\color{gray} Serendipitech Inc.}

\vfill

% Volume badge
\begin{tikzpicture}
  \node[
    draw=c2, rounded corners=4pt,
    inner sep=10pt,
    fill=c2!15,
    text=white
  ]{
    \ttfamily\small
    Vol.\,I \quad|\quad Foundations
  };
\end{tikzpicture}

\vspace{1cm}
{\small\color{gray}\today}

\end{center}
\end{titlepage}

% Restore normal page color
\pagecolor{white}\color{black}

% ── FRONT MATTER ────────────────────────────────────────────
\frontmatter
\tableofcontents

% ── Preface ─────────────────────────────────────────────────
\chapter*{Preface}
\addcontentsline{toc}{chapter}{Preface}

This monograph develops an interdisciplinary framework that imports concepts from
formal language theory and syntactic analysis into computational genomics.
The central thesis is that the genome is not merely a list of independent
molecular markers, but a structured object governed by hierarchical rules
analogous to the grammars of natural languages.

The intended audience is researchers trained in engineering, mathematics,
or computer science who wish to engage rigorously with genomic data.
Biological background is introduced from first principles wherever possible,
and each concept is grounded in formal notation familiar to the engineering
community.

\bigskip
\noindent
\textit{Chapter~1} establishes the physical and informational structure of
the genome, deriving a hierarchy from the primitive unit (the nucleotide
base) up to the full genome, expressed as a context-free generative grammar.

\vfill
\begin{flushright}
\textit{Wanhong}\\
Tokyo, \today
\end{flushright}

% ── MAIN MATTER ─────────────────────────────────────────────
\mainmatter

% ============================================================
\chapter{Genome Structure: A Hierarchical Grammar}
\label{chap:structure}
% ============================================================

\begin{notebox}[Chapter Overview]
This chapter introduces the physical architecture of the genome from the
bottom up.  We treat each structural level as a production rule in a
generative grammar, culminating in a six-level hierarchy that maps
naturally onto the formal language hierarchy of Chomsky.
\end{notebox}

\bigskip

% ────────────────────────────────────────────────────────────
\section{The Genome as a Formal Language}
\label{sec:genome-as-language}
% ────────────────────────────────────────────────────────────

A genome is the complete set of hereditary instructions of an organism,
encoded as a linear sequence over a four-letter alphabet
$\Sigma = \{A,\, T,\, G,\, C\}$.
From an engineering perspective this is a \emph{quaternary string} of length
approximately $3 \times 10^9$ characters in the case of \textit{Homo sapiens}.

\begin{definition}[Genome Alphabet]
Let $\Sigma = \{A, T, G, C\}$ be the terminal alphabet, where each symbol
represents one of the four nucleotide bases:
\begin{align*}
  A &\;=\; \text{Adenine},   &
  T &\;=\; \text{Thymine},   \\
  G &\;=\; \text{Guanine},   &
  C &\;=\; \text{Cytosine.}
\end{align*}
\end{definition}

Unlike a flat string, however, the genome possesses rich hierarchical
structure.  Bases group into codons, codons into exons, exons combine with
regulatory elements to form genes, genes are organised along chromosomes,
and chromosomes together constitute the genome.  This hierarchy is
reminiscent of the phrase-structure hierarchy in natural language,
and motivates a \emph{generative grammar} description.

% ────────────────────────────────────────────────────────────
\section{The Generative Grammar of Genomic Structure}
\label{sec:grammar}
% ────────────────────────────────────────────────────────────

We adopt the notation of context-free grammars (CFGs).
A production rule is written $\NT{A} \derives \alpha$, where
\NT{A} is a non-terminal (structural category) and $\alpha$ is a string
of non-terminals and/or terminals.  The operator ``$\gor$'' separates
alternatives; ``${}^+$'' means one or more occurrences; ``${}^*$''
means zero or more.

\begin{grammarbox}
\small
\renewcommand{\arraystretch}{1.5}
\begin{tabular}{@{}p{0.55\linewidth}@{\hspace{0.3em}}p{0.4\linewidth}@{}}
$\NT{Genome} \derives \NT{Chromosome}^{+}$ & \textcolor{commentcolor}{\scriptsize\textit{all chromosomes together}} \
$\NT{Chromosome} \derives \NT{Chromatid} \mathbin{{-}{-}} \NT{Chromatid}$ & \textcolor{commentcolor}{\scriptsize\textit{sister chromatid pair}} \
$\NT{Chromatid} \derives \NT{Locus}^{+}$ & \textcolor{commentcolor}{\scriptsize\textit{ordered sequence of loci}} \
$\NT{Locus} \derives \NT{Gene} \gor \NT{NonCoding}$ & \textcolor{commentcolor}{\scriptsize\textit{coding or non-coding region}} \
$\NT{NonCoding} \derives \NT{Base}^{+}$ & \textcolor{commentcolor}{\scriptsize\textit{non-coding DNA sequence}} \
\multicolumn{2}{@{}l}{\textbf{\scriptsize —— Core Gene Structure ——}} \
$\NT{Gene} \derives \NT{RegulatoryRegion};\NT{CodingRegion}$ & \textcolor{commentcolor}{\scriptsize\textit{regulatory $+$ coding}} \
$\NT{RegulatoryRegion} \derives \NT{Promoter};\NT{Enhancer}^{}$ & \textcolor{commentcolor}{\scriptsize\textit{promoter $+$ zero or more enhancers}} \
$\NT{Promoter} \derives \NT{CoreElement};\NT{TFBS}^{}$ & \textcolor{commentcolor}{\scriptsize\textit{core promoter $+$ proximal TF sites}} \
$\NT{CoreElement} \derives \T{T}\T{A}\T{T}\T{A}\T{A}\T{A} \gor \T{C}\T{A}\T{A}\T{T}$ & \textcolor{commentcolor}{\scriptsize\textit{TATA-box or CAAT-box}} \
$\NT{Enhancer} \derives \NT{TFBS}^{+}$ & \textcolor{commentcolor}{\scriptsize\textit{cluster of TF binding sites}} \
$\NT{TFBS} \derives \NT{Base};\NT{Base};\NT{Base};\NT{Base};\NT{Base};\NT{Base}$ & \textcolor{commentcolor}{\scriptsize\textit{minimal 6bp motif}} \
$\NT{CodingRegion} \derives \NT{Exon};(\NT{Intron};\NT{Exon})^{*}$ & \textcolor{commentcolor}{\scriptsize\textit{alternating exons and introns}} \
$\NT{Exon} \derives \NT{Codon}^{+}$ & \textcolor{commentcolor}{\scriptsize\textit{sequence of codons}} \
$\NT{Intron} \derives \NT{Base}^{+}$ & \textcolor{commentcolor}{\scriptsize\textit{non-coding intronic sequence}} \
$\NT{Codon} \derives \NT{Base};\NT{Base};\NT{Base}$ & \textcolor{commentcolor}{\scriptsize\textit{triplet of bases}} \
$\NT{Base} \derives \T{A} \gor \T{T} \gor \T{G} \gor \T{C}$ & \textcolor{commentcolor}{\scriptsize\textit{terminal symbols}} \
\end{tabular}
\smallskip\
\noindent
{\small\color{ntcolor}\textbf{Non-terminals}} are structural categories; {\small\color{tcolor}\textbf{Terminals}} $\T{A}, \T{T}, \T{G}, \T{C}$ are the irreducible atoms.
\end{grammarbox}

\begin{grammarbox}
\small
\renewcommand{\arraystretch}{1.5}
\begin{tabular}{@{}p{0.55\linewidth}@{\hspace{0.3em}}p{0.4\linewidth}@{}}
\multicolumn{2}{@{}l}{\textbf{\scriptsize —— Cell Cycle ——}} \
$\NT{CellCycle} \derives \NT{G1};\NT{S};\NT{G2};\NT{M}$ & \textcolor{commentcolor}{\scriptsize\textit{cell division cycle}} \
$\NT{G1} \derives \NT{G1Genes};\NT{Growth}$ & \textcolor{commentcolor}{\scriptsize\textit{cell growth phase}} \
$\NT{S} \derives \NT{ReplicationGenes};\NT{HistoneGenes}$ & \textcolor{commentcolor}{\scriptsize\textit{DNA synthesis phase}} \
$\NT{G2} \derives \NT{PreparationGenes}$ & \textcolor{commentcolor}{\scriptsize\textit{preparation for mitosis}} \
$\NT{M} \derives \NT{Division}$ & \textcolor{commentcolor}{\scriptsize\textit{mitosis phase}} \
\multicolumn{2}{@{}l}{\textbf{\scriptsize —— Phase-Specific Genes ——}} \
$\NT{G1Genes} \derives \NT{Gene}^{}$ & \textcolor{commentcolor}{\scriptsize\textit{genes active in G1}} \
$\NT{ReplicationGenes} \derives \NT{Gene}^{}$ & \textcolor{commentcolor}{\scriptsize\textit{DNA polymerase, etc.}} \
$\NT{HistoneGenes} \derives \NT{Gene}^{+}$ & \textcolor{commentcolor}{\scriptsize\textit{histone proteins}} \
$\NT{PreparationGenes} \derives \NT{Gene}^{*}$ & \textcolor{commentcolor}{\scriptsize\textit{tubulin, cyclins}} \
$\NT{Growth} \derives \epsilon$ & \textcolor{commentcolor}{\scriptsize\textit{cell growth process}} \
$\NT{Division} \derives \epsilon$ & \textcolor{commentcolor}{\scriptsize\textit{cell division process}} \
\end{tabular}
\smallskip\
\noindent
{\small\color{ntcolor}\textbf{Non-terminals}} are structural categories; {\small\color{tcolor}\textbf{Terminals}} $\T{A}, \T{T}, \T{G}, \T{C}$ are the irreducible atoms.
\end{grammarbox}

\begin{grammarbox}
\small
\renewcommand{\arraystretch}{1.5}
\begin{tabular}{@{}p{0.55\linewidth}@{\hspace{0.3em}}p{0.4\linewidth}@{}}
\multicolumn{2}{@{}l}{\textbf{\scriptsize —— Chromatin States ——}} \
$\NT{Locus} \derives \NT{Open} \gor \NT{Closed}$ & \textcolor{commentcolor}{\scriptsize\textit{chromatin accessibility}} \
$\NT{Open} \derives \NT{ActiveGene} \gor \NT{PoisedGene}$ & \textcolor{commentcolor}{\scriptsize\textit{euchromatin}} \
$\NT{Closed} \derives \NT{SilencedGene} \gor \NT{Repetitive}$ & \textcolor{commentcolor}{\scriptsize\textit{heterochromatin}} \
$\NT{ActiveGene} \derives \NT{Gene}$ & \textcolor{commentcolor}{\scriptsize\textit{actively transcribed}} \
$\NT{PoisedGene} \derives \NT{Gene}$ & \textcolor{commentcolor}{\scriptsize\textit{ready for activation}} \
$\NT{SilencedGene} \derives \NT{Gene}$ & \textcolor{commentcolor}{\scriptsize\textit{permanently repressed}} \
$\NT{Repetitive} \derives \NT{Base}^{+}$ & \textcolor{commentcolor}{\scriptsize\textit{repetitive DNA}} \
\multicolumn{2}{@{}l}{\textbf{\scriptsize —— RNA Processing ——}} \
$\NT{PrimaryTranscript} \derives \NT{Gene}$ & \textcolor{commentcolor}{\scriptsize\textit{initial transcript}} \
$\NT{MatureRNA} \derives \NT{Cap};\NT{Exon}^{+};\NT{Tail}$ & \textcolor{commentcolor}{\scriptsize\textit{processed mRNA}} \
$\NT{Cap} \derives \epsilon$ & \textcolor{commentcolor}{\scriptsize\textit{5' cap}} \
$\NT{Tail} \derives \T{A}\T{A}\T{A}\T{A}\T{A}\T{A}\T{A}\T{A}\T{A}\T{A}$ & \textcolor{commentcolor}{\scriptsize\textit{poly-A tail (simplified)}} \
\multicolumn{2}{@{}l}{\textbf{\scriptsize —— Protein ——}} \
$\NT{Protein} \derives \NT{AminoChain}$ & \textcolor{commentcolor}{\scriptsize\textit{functional protein}} \
$\NT{AminoChain} \derives \NT{Amino}^{+}$ & \textcolor{commentcolor}{\scriptsize\textit{polypeptide chain}} \
$\NT{Amino} \derives \NT{CodonToAmino}$ & \textcolor{commentcolor}{\scriptsize\textit{amino acid}} \
$\NT{CodonToAmino} \derives \T{Met} \gor \T{Phe} \gor \T{Leu} \gor \T{Ser} \gor \T{Pro} \gor \T{Arg}$ & \textcolor{commentcolor}{\scriptsize\textit{simplified mapping}} \
$\T{Met} \derives \epsilon$ \quad $\T{Phe} \derives \epsilon$ \quad $\T{Leu} \derives \epsilon$ & \textcolor{commentcolor}{\scriptsize\textit{amino acid symbols}} \
$\T{Ser} \derives \epsilon$ \quad $\T{Pro} \derives \epsilon$ \quad $\T{Arg} \derives \epsilon$ & \textcolor{commentcolor}{\scriptsize\textit{amino acid symbols}} \
\end{tabular}
\smallskip\
\noindent
{\small\color{ntcolor}\textbf{Non-terminals}} are structural categories; {\small\color{tcolor}\textbf{Terminals}} are $\T{A}, \T{T}, \T{G}, \T{C}$ and amino acid symbols. $\epsilon$ = empty production.
\end{grammarbox}

\begin{remark}
This grammar is deliberately simplified.  In reality, the production rules
for \NT{RegulatoryRegion} are context-sensitive (the same enhancer can
activate different genes depending on three-dimensional chromatin folding),
motivating extensions beyond pure CFGs in later chapters.
\end{remark}

% ────────────────────────────────────────────────────────────
\section{Level-by-Level Description of the Hierarchy}
\label{sec:levels}
% ────────────────────────────────────────────────────────────

The following subsections describe each level of the grammar,
progressing from the primitive terminal symbol upward to the genome.

% ────────────────────────────────────────────
\subsection{Level 1 — Nucleotide Base (Terminal Symbol)}
\label{subsec:base}
% ────────────────────────────────────────────

\begin{definition}[Nucleotide Base]
A \emph{nucleotide base} is one of the four organic molecules
$\{A, T, G, C\}$ that serve as the terminal symbols of the genomic alphabet.
Each base is separated from its neighbour by approximately $0.34\,\text{nm}$
along the DNA double helix.
\end{definition}

Bases pair according to \textbf{Watson--Crick complementarity}:
\begin{equation}
  A \;\Longleftrightarrow\; T \quad (\text{2 hydrogen bonds}),
  \qquad
  G \;\Longleftrightarrow\; C \quad (\text{3 hydrogen bonds}).
\end{equation}

This complementarity defines the double-stranded structure of DNA and is
fundamental to replication fidelity.  The two strands are antiparallel:

\begin{center}
\ttfamily\small
\begin{tabular}{ccccccccc}
5$'$ & A & T & G & C & G & A & T & \ldots \\
     & | & | & | & | & | & | & | & \\
3$'$ & T & A & C & G & C & T & A & \ldots
\end{tabular}
\end{center}

\begin{notebox}[Engineering Analogy]
Think of the base alphabet $\Sigma$ as the bit-set of a quaternary
information system.  Each base carries $\log_2 4 = 2\,\text{bits}$
of information.  The entire human haploid genome encodes approximately
$2 \times 3 \times 10^9 = 6 \times 10^9\,\text{bits} \approx 750\,\text{MB}$.
\end{notebox}

% ────────────────────────────────────────────
\subsection{Level 2 — Codon}
\label{subsec:codon}
% ────────────────────────────────────────────

\begin{definition}[Codon]
A \emph{codon} is a contiguous triplet of bases $b_1 b_2 b_3 \in \Sigma^3$.
It constitutes the minimal unit that specifies an amino acid or a
translational control signal.
\end{definition}

The codon space has cardinality $|\Sigma|^3 = 4^3 = 64$, yet only 20 amino
acids are encoded, together with one start codon (\texttt{ATG}) and three
stop codons (\texttt{TAA}, \texttt{TAG}, \texttt{TGA}).
The resulting many-to-one mapping is called the \emph{degeneracy} of the
genetic code:

\begin{equation}
  f : \Sigma^3 \longrightarrow \mathcal{A} \cup \{\text{start}, \text{stop}\},
  \quad |\mathcal{A}| = 20,
  \quad |{\Sigma^3}| = 64.
\end{equation}

\begin{keybox}[Engineering Analogy]
A codon is a 3-character opcode in a fixed-width instruction set.
The genetic code $f$ is the \emph{decode table} of this instruction set
architecture (ISA), mapping each 6-bit opcode to one of 21 operations
(20 amino acids + stop).
\end{keybox}

% ────────────────────────────────────────────
\subsection{Level 3 — Exon and Intron}
\label{subsec:exon-intron}
% ────────────────────────────────────────────

A gene's coding region is not a contiguous sequence of codons.
It is interrupted by non-coding segments:

\begin{definition}[Exon]
An \emph{exon} is a sequence of codons that is retained in the mature
messenger RNA (mRNA) after post-transcriptional processing.
Exons directly encode the amino acid sequence of a protein.
\end{definition}

\begin{definition}[Intron]
An \emph{intron} is an intervening sequence that is transcribed into
the precursor mRNA (pre-mRNA) but subsequently excised by the
\emph{spliceosome} machinery, leaving no trace in the mature mRNA.
\end{definition}

The linear structure of a pre-mRNA transcript is:
\begin{equation}
  5'\;\underbrace{E_1}_{\text{exon}}\;
     \underbrace{I_1}_{\text{intron}}\;
     \underbrace{E_2}_{\text{exon}}\;
     \underbrace{I_2}_{\text{intron}}\;
     \underbrace{E_3}_{\text{exon}}\;\ldots\;3'
\end{equation}

After splicing:
\begin{equation}
  \text{mature mRNA} = 5'\;E_1\;E_2\;E_3\;\ldots\;3'
\end{equation}

\begin{notebox}[Quantitative note]
In humans, introns account for approximately 90\% of the total length
of a typical gene, while exons are comparatively short.
Protein-coding exons together cover only $\approx 1.5\%$ of the entire genome.
\end{notebox}

\textbf{Alternative splicing.}
The spliceosome can join different subsets of exons, producing multiple
distinct mRNA isoforms from a single gene locus.  This is a form of
\emph{ambiguous parsing}: the same pre-mRNA sequence admits multiple
parse trees under different splicing grammars, each yielding a distinct
protein product.  This phenomenon will be revisited in Chapter~4.

% ────────────────────────────────────────────
\subsection{Level 4 — Gene}
\label{subsec:gene}
% ────────────────────────────────────────────

\begin{definition}[Gene]
A \emph{gene} is a functional unit of hereditary information consisting
of a \emph{regulatory region} and a \emph{coding region}, together
occupying a contiguous locus on a chromosome.
\end{definition}

The internal organisation of a canonical protein-coding gene is:

\begin{center}
\begin{tikzpicture}[
  x=1.0cm, y=0.7cm,
  every node/.style={font=\small},
  lbl/.style={font=\footnotesize\ttfamily, above=2pt}
]
  % Backbone
  \draw[thick, ->, >=Stealth] (-0.3,0) -- (11.8,0) node[right]{$3'$};
  \node at (-0.5,0) {$5'$};

  % Enhancer
  \fill[c6!25] (0.2,-0.25) rectangle (1.6,0.25);
  \draw[c6,  thick] (0.2,-0.25) rectangle (1.6,0.25);
  \node[lbl, c6] at (0.9,0.25) {Enhancer};

  % Promoter
  \fill[c4!25] (1.8,-0.25) rectangle (3.2,0.25);
  \draw[c4,  thick] (1.8,-0.25) rectangle (3.2,0.25);
  \node[lbl, c4] at (2.5,0.25) {Promoter};

  % TSS marker
  \draw[c4, dashed] (3.2,0.5) -- (3.2,-0.5);
  \node[font=\tiny, c4, below=2pt] at (3.2,-0.25) {TSS};

  % Exon 1
  \fill[c3!25] (3.4,-0.25) rectangle (4.6,0.25);
  \draw[c3, thick] (3.4,-0.25) rectangle (4.6,0.25);
  \node[lbl, c3] at (4.0,0.25) {Exon 1};

  % Intron 1
  \fill[gray!10] (4.6,-0.25) rectangle (6.2,0.25);
  \draw[gray,  thick] (4.6,-0.25) rectangle (6.2,0.25);
  \node[lbl, gray] at (5.4,0.25) {Intron 1};

  % Exon 2
  \fill[c3!25] (6.2,-0.25) rectangle (7.4,0.25);
  \draw[c3, thick] (6.2,-0.25) rectangle (7.4,0.25);
  \node[lbl, c3] at (6.8,0.25) {Exon 2};

  % Intron 2
  \fill[gray!10] (7.4,-0.25) rectangle (9.0,0.25);
  \draw[gray,  thick] (7.4,-0.25) rectangle (9.0,0.25);
  \node[lbl, gray] at (8.2,0.25) {Intron 2};

  % Exon 3
  \fill[c3!25] (9.0,-0.25) rectangle (10.2,0.25);
  \draw[c3, thick] (9.0,-0.25) rectangle (10.2,0.25);
  \node[lbl, c3] at (9.6,0.25) {Exon 3};

  % Terminator
  \fill[c5!20] (10.4,-0.25) rectangle (11.6,0.25);
  \draw[c5, thick] (10.4,-0.25) rectangle (11.6,0.25);
  \node[lbl, c5] at (11.0,0.25) {Terminator};

  % Braces
  \draw[decorate, decoration={brace, amplitude=5pt, mirror},
        c4, thick] (0.2,-0.45) -- (3.2,-0.45)
    node[midway, below=6pt, c4, font=\footnotesize]
    {Regulatory Region};
  \draw[decorate, decoration={brace, amplitude=5pt, mirror},
        c3, thick] (3.4,-0.45) -- (11.6,-0.45)
    node[midway, below=6pt, c3, font=\footnotesize]
    {Coding Region};

\end{tikzpicture}
\end{center}

\bigskip

\begin{definition}[Promoter]
A \emph{promoter} is a regulatory DNA sequence located immediately upstream
of a gene's transcription start site (TSS).  It is the binding platform for
RNA Polymerase~II and associated general transcription factors, and is
\emph{required} for any transcription to occur.  In the grammar,
$\NT{Promoter} \derives \NT{CoreElement}\;\NT{TFBS}^{*}$, where
\NT{CoreElement} typically contains the \textbf{TATA box} (consensus
\texttt{TATAAA}, $\approx$30\,bp upstream of TSS) and/or Initiator element.
\end{definition}

\begin{definition}[Enhancer]
An \emph{enhancer} is a \emph{cis}-regulatory DNA element that amplifies
transcription from a target promoter by 10- to 1000-fold.
Unlike promoters, enhancers can be located \emph{anywhere} relative to their
target gene: upstream, downstream, or within an intron, at distances up to
$10^6$\,bp away, acting through DNA looping.  In the grammar,
$\NT{Enhancer} \derives \NT{TFBS}^{+}$: an enhancer is defined by a cluster
of transcription factor binding sites that act cooperatively.
\end{definition}

\begin{definition}[Transcription Factor Binding Site (TFBS) and Core Element]
A \emph{transcription factor binding site} (\NT{TFBS}) is a short DNA
sequence motif (typically 6--20\,bp) recognised and bound by a specific
transcription factor (TF) protein.  $\NT{TFBS} \derives \NT{Base}^{+}$.
A \emph{core element} (\NT{CoreElement}) is the minimal promoter-specific
sequence required for RNA Pol~II positioning, such as the TATA box or
CpG island; formally $\NT{CoreElement} \derives \NT{Base}^{+}$.
\end{definition}

\begin{definition}[Allele]
An \emph{allele} is one of two or more alternative versions of a gene
(or more generally, a genomic locus), differing in their nucleotide
sequence.  In a diploid organism, each locus is represented by two
alleles, one inherited from each parent.
\end{definition}

\begin{notebox}[Quantitative note]
The human genome contains approximately $20{,}000$ protein-coding genes,
which together span only $\approx 1.5\%$ of the 3.2\,Gbp genome.
The remaining $98.5\%$ consists of introns, regulatory elements,
transposable elements, pseudogenes, and sequences of unknown function.
\end{notebox}

% ────────────────────────────────────────────
\subsection{Level 5 — Chromosome}
\label{subsec:chromosome}
% ────────────────────────────────────────────

\begin{definition}[Chromosome]
A \emph{chromosome} is a single, continuous DNA molecule complexed with
histone proteins to form \emph{chromatin}.  During cell division it is
further condensed into a compact, rod-shaped structure visible by light
microscopy.
\end{definition}

The physical compaction hierarchy is:

\begin{equation}
  \text{DNA double helix}
  \;\xrightarrow{\;\text{wrap}\;}\;
  \text{nucleosome}
  \;\xrightarrow{\;\text{fold}\;}\;
  \text{chromatin fibre}
  \;\xrightarrow{\;\text{scaffold}\;}\;
  \text{chromosome}
\end{equation}

This $\sim\!10{,}000\times$ linear compaction allows $\approx 2\,\text{m}$
of total DNA to fit within a nucleus of diameter $\approx 6\,\mu\text{m}$.

\textbf{Human karyotype.}
Humans possess 46 chromosomes arranged in 23 pairs: 22 pairs of
\emph{autosomes} (Chr1--Chr22, largest to smallest) plus one pair of
\emph{sex chromosomes} (XX in females, XY in males).

\begin{center}
\begin{tikzpicture}[x=0.55cm, y=0.5cm]
  \foreach \i in {1,...,22}{
    \pgfmathsetmacro{\h}{3.5 - 0.12*\i}
    \pgfmathtruncatemacro{\x}{(\i-1)}
    \fill[c5!60, rounded corners=2pt]
      (\x, 0) rectangle (\x+0.35, \h);
    \fill[c5!40, rounded corners=2pt]
      (\x+0.4, 0) rectangle (\x+0.75, \h*0.97);
    \node[font=\tiny, rotate=90] at (\x+0.37, -0.35) {\i};
  }
  % Sex chromosomes
  \fill[c2!60, rounded corners=2pt] (22.2, 0) rectangle (22.55, 3.0);
  \fill[c2!40, rounded corners=2pt] (22.6, 0) rectangle (22.95, 3.0);
  \node[font=\tiny] at (22.57, -0.35) {XX};
  \fill[c1!60, rounded corners=2pt] (23.3, 0) rectangle (23.65, 3.0);
  \fill[c1!30, rounded corners=2pt] (23.7, 0) rectangle (24.05, 1.2);
  \node[font=\tiny] at (23.67, -0.35) {XY};
  \node[font=\scriptsize, gray] at (12, -0.9) {Chromosome pairs 1--22 (autosomes) + sex chromosomes};
\end{tikzpicture}
\end{center}

\begin{definition}[Locus (plural: Loci)]
A \emph{locus} is a fixed, named position on a chromosome where a
particular gene or other genomic feature resides.
Loci are specified by chromosome number, arm (p = short, q = long),
and band coordinates, e.g.\ \texttt{17q21.31}.
\end{definition}

% ────────────────────────────────────────────
\subsection{Level 6 — Genome}
\label{subsec:genome}
% ────────────────────────────────────────────

\begin{definition}[Genome]
The \emph{genome} of an organism is the complete set of its genetic
information, comprising all chromosomes.  For diploid species the genome
contains two copies of each autosome plus two sex chromosomes.
The \emph{haploid} genome refers to a single copy of each chromosome.
\end{definition}

\begin{table}[h]
\centering
\caption{Quantitative summary of the human genome.}
\label{tab:genome-stats}
\begin{tabularx}{\textwidth}{lXr}
\toprule
\textbf{Feature} & \textbf{Description} & \textbf{Value} \\
\midrule
Haploid size         & Total base pairs (haploid set)       & $3.2\,\Gbp$ \\
Diploid size         & Both homologous sets                 & $6.4\,\Gbp$ \\
Chromosomes          & Pairs in somatic cells               & 23 pairs (46) \\
Protein-coding genes & Distinct gene loci                   & $\approx 20{,}000$ \\
Coding fraction      & Fraction of genome in exons          & $\approx 1.5\%$ \\
Regulatory fraction  & Estimated functional non-coding       & $\approx 8\%$ \\
Known SNPs           & Common single-nucleotide variants    & $\approx 10^7$ \\
Information content  & At 2 bits/base (haploid)             & $\approx 750\,\text{MB}$ \\
\bottomrule
\end{tabularx}
\end{table}

% ────────────────────────────────────────────────────────────
\section{The Full Hierarchy Diagram}
\label{sec:hierarchy-diagram}
% ────────────────────────────────────────────────────────────

Figure~\ref{fig:hierarchy} synthesises the six levels into a single
parse-tree style diagram.  Each node is annotated with its grammar category,
physical scale, and engineering analogue.

\begin{figure}[p]
\centering
\begin{tikzpicture}[
  x=1cm, y=1cm,
  node distance=0pt,
  level/.style={
    draw, rounded corners=5pt,
    minimum width=13cm, text width=12cm, minimum height=1.5cm,
    inner sep=10pt, align=left,
    font=\small
  },
  lnum/.style={
    font=\footnotesize\bfseries\ttfamily,
    text=white, fill=#1,
    minimum width=1.2cm, minimum height=1.5cm,
    inner sep=4pt, align=center
  },
  arr/.style={->, >=Stealth, thick, #1},
  alabel/.style={font=\footnotesize\itshape, fill=white, inner sep=2pt}
]

% ── Level 1: Base ──────────────────────────────────────────
\node[level, fill=c1!10, draw=c1]
  (L1) at (0.6, 0)
  {\textbf{\textcolor{c1}{Base}} \enspace
   $\NT{Base} \derives \T{A} \gor \T{T} \gor \T{G} \gor \T{C}$
   \hspace{1em} \textcolor{gray}{\footnotesize 0.34\,nm / base}\\[3pt]
   \textit{Terminal symbol.  4-ary digit.
   Watson--Crick pairing: A$=$T (2H), G$\equiv$C (3H).}};
\node[lnum=c1] at (-6.0, 0) {L1};

% ── Level 2: Codon ─────────────────────────────────────────
\node[level, fill=c2!8, draw=c2]
  (L2) at (0.6,-2.6)
  {\textbf{\textcolor{c2}{Codon}} \enspace
   $\NT{Codon} \derives \NT{Base}\;\NT{Base}\;\NT{Base}$
   \hspace{1em} \textcolor{gray}{\footnotesize 3\,bp}\\[3pt]
   \textit{64 possible codons $\to$ 20 amino acids + start/stop.
   Analogous to a fixed-width ISA opcode.}};
\node[lnum=c2] at (-6.0,-2.6) {L2};

% ── Level 3: Exon/Intron ────────────────────────────────────
\node[level, fill=c3!8, draw=c3]
  (L3) at (0.6,-5.2)
  {\textbf{\textcolor{c3}{Exon / Intron}} \enspace
   $\NT{Exon}\!\derives\!\NT{Codon}^{+};\;
    \NT{CodingRegion}\!\derives\!\NT{Exon}(\NT{Intron}\;\NT{Exon})^{*}$
   \hspace{0.5em}\textcolor{gray}{\footnotesize 100\,bp--100\,kbp}\\[3pt]
   \textit{Exon: retained in mature mRNA (meaningful words).
   Intron: excised by spliceosome (compiler-discarded comments).}};
\node[lnum=c3] at (-6.0,-5.2) {L3};

% ── Level 4: Gene ──────────────────────────────────────────
\node[level, fill=c4!8, draw=c4]
  (L4) at (0.6,-7.8)
  {\textbf{\textcolor{c4}{Gene}} \enspace
   $\NT{Gene} \derives \NT{RegulatoryRegion}\;\NT{CodingRegion}$
   \hspace{1em} \textcolor{gray}{\footnotesize 1\,kbp--2\,Mbp}\\[3pt]
   \textit{Functional unit encoding one protein (or RNA).
   $\approx$20\,000 in humans; only 1.5\% of genome.
   Regulatory region expands to: $\NT{Promoter}\;\NT{Enhancer}^*$.}};
\node[lnum=c4] at (-6.0,-7.8) {L4};

% ── Level 5: Chromosome ────────────────────────────────────
\node[level, fill=c5!8, draw=c5]
  (L5) at (0.6,-10.4)
  {\textbf{\textcolor{c5}{Chromosome}} \enspace
   $\NT{Chromosome} \derives \NT{Chromatid}\mathbin{{-}{-}}\NT{Chromatid}$;\enspace
   $\NT{Chromatid} \derives \NT{Locus}^{+}$
   \hspace{1em}\textcolor{gray}{\footnotesize 50--250\,Mbp}\\[3pt]
   \textit{Continuous DNA + histone scaffold.
   Humans: 46 chromosomes (23 pairs).
   Analogous to a source-code module or file.}};
\node[lnum=c5] at (-6.0,-10.4) {L5};

% ── Level 6: Genome ────────────────────────────────────────
\node[level, fill=c6!8, draw=c6]
  (L6) at (0.6,-13.0)
  {\textbf{\textcolor{c6}{Genome}} \enspace
   $\NT{Genome} \derives \NT{Chromosome}^{+}$
   \hspace{1em} \textcolor{gray}{\footnotesize $\approx$3.2\,Gbp (haploid)}\\[3pt]
   \textit{Complete hereditary information of the organism.
   $\approx$750\,MB at 2\,bits/base.
   Analogous to the full codebase / operating system.}};
\node[lnum=c6] at (-6.0,-13.0) {L6};

% ── Arrows ─────────────────────────────────────────────────
\draw[arr=c2]  (L1.south) -- node[alabel]{$\times 3$} (L2.north);
\draw[arr=c3]  (L2.south) -- node[alabel]{sequences of codons} (L3.north);
\draw[arr=c4]  (L3.south) -- node[alabel]{regulatory + coding regions} (L4.north);
\draw[arr=c5]  (L4.south) -- node[alabel]{many genes + non-coding regions} (L5.north);
\draw[arr=c6]  (L5.south) -- node[alabel]{23 pairs} (L6.north);

\end{tikzpicture}
\caption{Six-level generative hierarchy of genomic structure.
  Each row corresponds to one production rule in the genome grammar
  (\S\ref{sec:grammar}).
  Left badges give the level number; right annotations give the physical
  scale.  Arrows are labelled with the composition relationship.}
\label{fig:hierarchy}
\end{figure}

% ────────────────────────────────────────────────────────────
\section{Comparative Summary: Gene, Chromosome, Genome}
\label{sec:comparison}
% ────────────────────────────────────────────────────────────

A common source of confusion is the distinction between the three terms
that name the central objects of this monograph.
Table~\ref{tab:compare} provides a side-by-side comparison.

\begin{table}[htbp]
\centering
\caption{Comparison of the three primary structural concepts.}
\label{tab:compare}
\renewcommand{\arraystretch}{1.4}
{\small
\begin{tabularx}{\textwidth}{>{\bfseries}l X X X}
\toprule
Concept & Definition & Scale & Engineering analogue \\
\midrule
\textcolor{c4}{Gene}
  & Encodes one protein or functional RNA; includes
    regulatory and coding sub-regions
  & 1\,kbp--2\,Mbp; $\approx$20\,000 in humans
  & Function / subroutine \\
\textcolor{c5}{Chromosome}
  & Continuous DNA + histone scaffold; carries many genes
    and non-coding sequence
  & 50--250\,Mbp; 46 in humans
  & Source-code module \\
\textcolor{c6}{Genome}
  & Complete hereditary information; sum of all chromosomes
  & $\approx$3.2\,Gbp haploid; $\approx$750\,MB
  & Full codebase / OS \\
\midrule
\textcolor{c2}{Locus}
  & Named position on a chromosome; addressing mechanism,
    not an independent physical entity
  & Precise address
  & Memory address \\
\textcolor{c1}{Allele}
  & Alternative sequence at the same locus; differs between
    individuals or homologous chromosomes
  & Gene-sized
  & Alternative function version \\
\bottomrule
\end{tabularx}
}
\end{table}

\begin{keybox}[Critical Insight for Chapters Ahead]
All cells of an organism contain \emph{identical} genome sequences
(with rare exceptions).  Yet a liver cell and a neuron are functionally
and morphologically distinct.  This difference arises entirely from
\emph{which genes are expressed}, and at what level --- a fact controlled
by the regulatory regions (promoters and enhancers) described in
\S\ref{subsec:gene}.  The grammar rules governing these regulatory
elements are the central target of genomic syntactic analysis.
\end{keybox}

% ────────────────────────────────────────────────────────────
\section{Chapter Summary}
\label{sec:ch1-summary}
% ────────────────────────────────────────────────────────────

This chapter established the following facts, which will be used
throughout the remainder of the monograph:

\begin{enumerate}
  \item The genome is a quaternary string over $\Sigma = \{A,T,G,C\}$,
        carrying $\approx 750\,\text{MB}$ of information.
  \item Its structure is hierarchical and can be expressed as a
        context-free generative grammar with six levels: Base $\to$
        Codon $\to$ Exon/Intron $\to$ Gene $\to$ Chromosome $\to$ Genome.
  \item A \emph{gene} ($\approx 20{,}000$ in humans) is a functional
        unit comprising a regulatory region and a coding region.
  \item A \emph{chromosome} is the physical carrier of many genes plus
        extensive non-coding sequence.
  \item The \emph{genome} is the union of all chromosomes and represents
        the complete hereditary specification.
  \item Only $\approx 1.5\%$ of the genome encodes proteins; the
        remaining $98.5\%$ includes regulatory elements, introns,
        repetitive sequences, and regions of currently unknown function.
        This ``dark matter'' of the genome is where syntactic analysis
        offers the most analytical leverage.
\end{enumerate}

\noindent
Chapter~2 introduces the principal tool for detecting genetic variation
--- Genome-Wide Association Studies (GWAS) --- and analyses its
statistical foundations and limitations from the perspective of a
researcher trained in engineering and formal methods.

% ── End of Chapter 1 ────────────────────────────────────────
\end{document}